\subsection*{Motivation}

Chapter~28 framed RLHF as KL-regularized reward maximization over the entire response distribution and derived DPO and PPO at the \emph{loss} level. To understand \emph{why} those algorithms converge --- and to make sense of value baselines, advantages, and the policy gradient itself --- we need the formal object underneath every reinforcement-learning method: the \emph{Markov Decision Process} (MDP). This chapter builds MDPs from scratch, derives the Bellman expectation and optimality equations as direct consequences of the definition of expected return (Chapter~10), proves the Bellman operator is a $\gamma$-contraction in sup norm, invokes the Banach fixed-point theorem of Chapter~7 to conclude that value iteration converges geometrically, sketches finite-step convergence of policy iteration, and quotes the Watkins--Tsitsiklis theorem connecting tabular Q-learning to the Robbins--Monro stochastic-approximation theory of Chapter~13. Token-by-token generation by an LLM is itself an MDP --- state $=$ prefix, action $=$ next token, reward $=$ terminal preference signal --- so every theorem here applies, modulo the combinatorial blow-up of $|\mathcal{S}|$, to the policies trained in Chapters 30--31.

\subsection*{Definitions}

\begin{definition}[Markov Decision Process]
A tuple $(\mathcal{S}, \mathcal{A}, P, r, \gamma)$ with $\mathcal{S}$ a finite state space, $\mathcal{A}$ a finite action space, $P : \mathcal{S}\times\mathcal{A} \to \Delta(\mathcal{S})$ a transition kernel ($P(s'\mid s,a)$ is the conditional probability of $s'$ given $(s,a)$, in the sense of Chapter~8), $r : \mathcal{S}\times\mathcal{A} \to \mathbb{R}$ a bounded reward, and $\gamma \in [0,1)$ a discount factor.
\end{definition}

\begin{definition}[Trajectory and return]
A sample path is $\tau = (s_0, a_0, r_0, s_1, a_1, r_1, \dots)$ generated by $a_t \sim \pi(\cdot\mid s_t)$, $s_{t+1}\sim P(\cdot\mid s_t, a_t)$, $r_t = r(s_t, a_t)$. The \emph{discounted return} is
\[
G_t \;:=\; \sum_{k=0}^{\infty}\gamma^k r_{t+k}.
\]
Since $|r|\le R_{\max}$ and $\gamma<1$, $|G_t|\le R_{\max}/(1-\gamma)$ and the series converges absolutely.
\end{definition}

\begin{definition}[Stationary policy and value functions]
A stationary policy is $\pi : \mathcal{S}\to\Delta(\mathcal{A})$. Define
\[
V^\pi(s) := \mathbb{E}_\pi[G_t\mid s_t=s], \qquad
Q^\pi(s,a) := \mathbb{E}_\pi[G_t\mid s_t=s, a_t=a],
\]
\[
V^*(s) := \sup_\pi V^\pi(s), \qquad Q^*(s,a) := \sup_\pi Q^\pi(s,a).
\]
\end{definition}

\begin{definition}[Bellman operators]
For $V \in \mathbb{R}^{|\mathcal{S}|}$,
\[
(\mathcal{T}^\pi V)(s) := \sum_a \pi(a\mid s)\!\left[r(s,a) + \gamma \sum_{s'} P(s'\mid s,a)\, V(s')\right],
\]
\[
(\mathcal{T}^* V)(s) := \max_a\!\left[r(s,a) + \gamma \sum_{s'} P(s'\mid s,a)\, V(s')\right].
\]
\end{definition}

\subsection*{Theorems and proofs}

\begin{theorem}[Bellman expectation equation]
$V^\pi = \mathcal{T}^\pi V^\pi$.
\end{theorem}

\begin{proof}
By definition $G_t = r_t + \gamma G_{t+1}$. Take $\mathbb{E}_\pi[\,\cdot\mid s_t = s]$ and apply linearity of expectation (Chapter~10):
\[
V^\pi(s) = \mathbb{E}_\pi[r_t \mid s_t=s] + \gamma \,\mathbb{E}_\pi[G_{t+1}\mid s_t=s].
\]
By the tower property,
\[
\mathbb{E}_\pi[r_t\mid s_t=s] = \sum_a \pi(a\mid s)\, r(s,a),
\]
and using the Markov property of $P$,
\[
\mathbb{E}_\pi[G_{t+1}\mid s_t=s] = \sum_{a,s'} \pi(a\mid s)\, P(s'\mid s,a)\, V^\pi(s').
\]
Combining gives $V^\pi(s) = (\mathcal{T}^\pi V^\pi)(s)$.
\end{proof}

\begin{theorem}[Bellman optimality equation]
$V^* = \mathcal{T}^* V^*$, and there exists a deterministic policy $\pi^*$ attaining the maximum.
\end{theorem}

\begin{proof}
Apply the same calculation at the optimal policy. The Bellman principle of optimality says that, conditional on being in state $s$, the optimal expected return decomposes as the maximum over the next action $a$ of the immediate reward $r(s,a)$ plus the discounted optimal expected return from $s' \sim P(\cdot\mid s,a)$, giving $V^*(s) = \max_a [r(s,a) + \gamma\sum_{s'} P(s'\mid s,a) V^*(s')]$. Because $\mathcal{A}$ is finite, the $\arg\max$ exists; defining $\pi^*(s)\in\arg\max$ produces a deterministic optimal policy.
\end{proof}

\begin{theorem}[Contraction]
$\mathcal{T}^\pi$ and $\mathcal{T}^*$ are $\gamma$-contractions in sup norm $\|V\|_\infty := \max_s |V(s)|$.
\end{theorem}

\begin{proof}
For any $V, V' \in \mathbb{R}^{|\mathcal{S}|}$ and any state $s$,
\[
|(\mathcal{T}^\pi V)(s) - (\mathcal{T}^\pi V')(s)|
= \gamma\left| \sum_a \pi(a\mid s) \sum_{s'} P(s'\mid s,a)\, [V(s')-V'(s')] \right|.
\]
Bounding the absolute value by the absolute sum and using $\sum_a \pi(a\mid s) = \sum_{s'} P(s'\mid s,a) = 1$ yields
\[
|(\mathcal{T}^\pi V)(s) - (\mathcal{T}^\pi V')(s)| \le \gamma \|V - V'\|_\infty.
\]
Taking the max over $s$ gives the sup-norm bound. For $\mathcal{T}^*$, use the elementary inequality $|\max_a f(a) - \max_a g(a)| \le \max_a |f(a) - g(a)|$ to reduce to the previous case.
\end{proof}

\begin{corollary}[Value iteration converges geometrically]
Let $V_{k+1} := \mathcal{T}^* V_k$. By Banach's fixed-point theorem (Chapter~7) on the complete normed space $(\mathbb{R}^{|\mathcal{S}|},\|\cdot\|_\infty)$,
\[
\|V_k - V^*\|_\infty \;\le\; \gamma^k \|V_0 - V^*\|_\infty.
\]
In particular, $\log\|V_k - V^*\|_\infty$ has slope $\log \gamma$ as a function of $k$.
\end{corollary}

\begin{theorem}[Policy iteration; finite termination]
Alternate (i) policy evaluation $V^\pi = (I - \gamma P^\pi)^{-1} r^\pi$, where $P^\pi(s,s') = \sum_a \pi(a\mid s)P(s'\mid s,a)$; the matrix $I - \gamma P^\pi$ is invertible because $\rho(\gamma P^\pi)\le \gamma <1$. (ii) Greedy improvement $\pi'(s) := \arg\max_a [r(s,a) + \gamma \sum_{s'} P(s'\mid s,a) V^\pi(s')]$. Then $V^{\pi'}\ge V^\pi$ pointwise, and the algorithm terminates in finitely many iterations at $\pi^*$.
\end{theorem}

\begin{proof}[Sketch]
By construction, $\mathcal{T}^{\pi'} V^\pi \ge V^\pi = \mathcal{T}^\pi V^\pi$. Applying $\mathcal{T}^{\pi'}$ repeatedly to both sides and using monotonicity ($V\le V'\Rightarrow \mathcal{T}^{\pi'} V \le \mathcal{T}^{\pi'} V'$) plus the contraction limit $(\mathcal{T}^{\pi'})^k V^\pi \to V^{\pi'}$ yields $V^{\pi'} \ge V^\pi$. The deterministic policy space has size $|\mathcal{A}|^{|\mathcal{S}|} <\infty$; each strict improvement visits a new policy; once $\pi' = \pi$, $\pi$ is a fixed point of $\mathcal{T}^*$, hence equals $\pi^*$.
\end{proof}

\begin{theorem}[Q-learning convergence; Watkins 1989, Tsitsiklis 1994]
Consider the update
\[
Q_{t+1}(s_t,a_t) = (1-\alpha_t)\,Q_t(s_t,a_t) + \alpha_t\!\left[r_t + \gamma\max_{a'} Q_t(s_{t+1}, a')\right].
\]
If every $(s,a)$ is visited infinitely often and the per-pair step sizes obey the Robbins--Monro conditions of Chapter~13, $\sum_t \alpha_t = \infty$, $\sum_t \alpha_t^2 <\infty$, then $Q_t \to Q^*$ almost surely.
\end{theorem}

\begin{proof}[Sketch]
The update is precisely the Robbins--Monro stochastic-approximation recursion for the fixed point of $\mathcal{T}^*$, which is a $\gamma$-contraction in sup norm. The sample $s_{t+1}\sim P(\cdot\mid s_t,a_t)$ replaces the expectation $\sum_{s'} P(s'\mid s,a)$ by an unbiased one-sample estimator, producing a martingale-difference noise with bounded variance (since rewards and $Q$ are bounded). Tsitsiklis' generalization of stochastic approximation to contraction operators gives almost-sure convergence under the stated step-size and visitation conditions.
\end{proof}

\begin{lemma}[Performance-difference; Kakade \& Langford, 2002]
\label{lem:pdl}
For any two policies $\pi, \pi'$ on the same MDP and any state $s_0$,
\[
  J(\pi') - J(\pi) \;=\; \frac{1}{1 - \gamma}\, \mathbb{E}_{(s, a) \sim d^{\pi'}, \pi'(\cdot \mid s)}\!\left[ A^\pi(s, a) \right],
\]
where $A^\pi(s, a) = Q^\pi(s, a) - V^\pi(s)$ is the advantage under $\pi$ and $d^{\pi'}(s) := (1 - \gamma) \sum_{t=0}^\infty \gamma^t \, \mathbb{P}^{\pi'}(s_t = s \mid s_0)$ is the discounted state-visitation distribution under $\pi'$.
\end{lemma}

\begin{proof}
By definition, $J(\pi') = \mathbb{E}^{\pi'}\!\left[\sum_{t=0}^\infty \gamma^t r_t\right]$, where $\mathbb{E}^{\pi'}$ denotes expectation under trajectories generated by $\pi'$ starting from $s_0$. The proof is a telescoping identity that inserts $V^\pi$ at every step. For any state $s$, the Bellman expectation equation for $V^\pi$ states
\[
  V^\pi(s) = \mathbb{E}_{a \sim \pi(\cdot \mid s)}\!\left[ r(s, a) + \gamma\, \mathbb{E}_{s' \sim P(\cdot \mid s, a)} V^\pi(s') \right].
\]
We DO NOT use this directly; instead, observe the trivial identity
\[
  J(\pi) = V^\pi(s_0) = \mathbb{E}^{\pi'}[V^\pi(s_0)] = \mathbb{E}^{\pi'}\!\left[\,V^\pi(s_0) + \sum_{t=0}^\infty \gamma^t \bigl(\gamma\, V^\pi(s_{t+1}) - V^\pi(s_{t+1})\bigr)\right],
\]
where we appended a telescoping zero (each $V^\pi(s_{t+1})$ term cancels its scaled copy at step $t+1$, modulo a $\lim_{T\to\infty} \gamma^{T+1} V^\pi(s_{T+1}) = 0$ tail that vanishes since $V^\pi$ is bounded and $\gamma < 1$). Rearranging,
\[
  J(\pi) = \mathbb{E}^{\pi'}\!\left[\sum_{t=0}^\infty \gamma^t \bigl(V^\pi(s_t) - \gamma\, V^\pi(s_{t+1})\bigr) - V^\pi(s_0)\right] + V^\pi(s_0) = \mathbb{E}^{\pi'}\!\left[\sum_{t=0}^\infty \gamma^t \bigl(V^\pi(s_t) - \gamma V^\pi(s_{t+1})\bigr)\right].
\]
Subtracting from $J(\pi')$,
\[
  J(\pi') - J(\pi) = \mathbb{E}^{\pi'}\!\left[\sum_{t=0}^\infty \gamma^t \bigl(r_t + \gamma V^\pi(s_{t+1}) - V^\pi(s_t)\bigr)\right].
\]
Inside the expectation, condition on $(s_t, a_t)$: the inner expression is exactly $A^\pi(s_t, a_t)$, since
\[
  \mathbb{E}\!\left[r_t + \gamma V^\pi(s_{t+1}) \mid s_t, a_t\right] = Q^\pi(s_t, a_t),
\]
and $V^\pi(s_t)$ is $\sigma(s_t)$-measurable. Hence
\[
  J(\pi') - J(\pi) = \mathbb{E}^{\pi'}\!\left[\sum_{t=0}^\infty \gamma^t A^\pi(s_t, a_t)\right] = \frac{1}{1-\gamma} \, \mathbb{E}_{(s,a) \sim d^{\pi'}, \pi'}\!\left[A^\pi(s, a)\right],
\]
where the last equality folds the discount factor into the discounted state-visitation $d^{\pi'}$. \qedhere
\end{proof}

\begin{remark}
The performance-difference lemma is the foundation of every \emph{conservative policy iteration} algorithm (Kakade \& Langford 2002), and via the next step --- replacing $d^{\pi'}$ by the on-policy $d^\pi$ and absorbing the substitution error into a KL penalty --- underlies TRPO and PPO (Chapter~31). It says exactly: the gain from switching $\pi \to \pi'$ is the expected $\pi$-advantage of $\pi'$'s actions, weighted by where $\pi'$ \emph{actually visits}.
\end{remark}

\subsection*{Code sketch}

The notebook builds a $4{\times}4$ GridWorld with reward $-1$ per step and $+10$ at an absorbing goal corner, runs value iteration with $\gamma = 0.9$, and verifies that the slope of $\log \|V_k - V^*\|_\infty$ versus $k$ matches $\log\gamma$. It extracts the greedy policy from $Q^*$, re-derives $V^*$ via policy iteration in dramatically fewer outer iterations, then runs tabular Q-learning for 5000 episodes with $\alpha_t = 1/(1+\mathrm{visits}(s,a))$ and $\varepsilon$-greedy exploration ($\varepsilon=0.1$), confirming $\|Q_t - Q^*\|_\infty \to 0$. A final cell runs REINFORCE on a 3-state, 2-action MDP with softmax policy parameterization and Monte-Carlo returns, previewing Chapter~30.

\subsection*{Connection to LLMs / RLHF}

A language model is an MDP whose state is the prompt-plus-prefix $s_t = (x, y_{<t})$, action is the next token $a_t = y_t \in \mathcal{V}$, transition is deterministic ($s_{t+1} = s_t \oplus a_t$), and the reward is sparse and terminal: $r_t = 0$ for $t<T$ and $r_T = r_\phi(x,y)$ from a learned reward model. Under this MDP the KL-regularized objective of Chapter~28 is exactly expected discounted return with per-step reward shaped by $-\beta \log\frac{\pi_\theta(a_t\mid s_t)}{\pi_{\mathrm{ref}}(a_t\mid s_t)}$. The Bellman optimality equation in this terminal-reward MDP collapses to the closed-form
\[
\pi^*(y\mid x) \propto \pi_{\mathrm{ref}}(y\mid x)\exp(r(x,y)/\beta),
\]
which is precisely the central DPO identity (Chapter~28, Theorem~1). Value iteration and policy iteration are intractable because $|\mathcal{S}|$ is exponential in context length, forcing parametric methods --- function approximation plus policy gradients --- the subject of Chapter~30 (REINFORCE, advantage estimation, GAE) and Chapter~31 (actor-critic / PPO machinery used in InstructGPT, Claude, and Llama-3).
